%%%%%%%%%%%%%%%%%%%%%%%%%%%%%%%%%%%%%%%%%
% Important note:
% This template must be compiled with XeLaTeX, the below lines will ensure this
%!TEX TS-program = xelatex
%!TEX encoding = UTF-8 Unicode
%
%%%%%%%%%%%%%%%%%%%%%%%%%%%%%%%%%%%%%%%%%

%----------------------------------------------------------------------------------------
%	PACKAGES AND OTHER DOCUMENT CONFIGURATIONS
%----------------------------------------------------------------------------------------

\documentclass[11pt, letterpaper]{awesome-cv} % A4 paper size by default, use 'letterpaper' for US letter

\geometry{left=2cm, top=1.5cm, right=2cm, bottom=2cm, footskip=.5cm} % Configure page margins with geometry
 
\fontdir[fonts/] % Specify the location of the included fonts

% Color for highlights
%\colorlet{awesome}{awesome-skyblue} % Default colors include: awesome-emerald, awesome-skyblue, awesome-red, awesome-pink, awesome-orange, awesome-nephritis, awesome-concrete, awesome-darknight
%\definecolor{awesome}{HTML}{CA63A8} % Uncomment if you would like to specify your own color
\definecolor{awesome}{HTML}{0BA89A} % UCSB sea green

% Colors for text - uncomment and modify
%\definecolor{darktext}{HTML}{414141}
%\definecolor{text}{HTML}{414141}
%\definecolor{graytext}{HTML}{414141}
%\definecolor{lighttext}{HTML}{414141}

\renewcommand{\acvHeaderSocialSep}{\quad\textbar\quad} % If you would like to change the social information separator from a pipe (|) to something else

%----------------------------------------------------------------------------------------
%	PERSONAL INFORMATION
%	Comment any of the lines below if they are not required
%----------------------------------------------------------------------------------------

\name{Jason C.}{Wei}
\address{2320 Willet Way, Pleasanton, California, 94566, USA}
\mobile{(925) 352-7600}
\github{github.com/Jayweiii}
\homepage{jweii.com}
\email{jwei@ucsb.edu}
\linkedin{linkedin.com/in/jwei}
%\skype{skypeid}
%\stackoverflow{SOid}{SOname}
%\twitter{@twit}
%\linkedin{linkedin name}
%\reddit{reddit account}
%\xing{xing name}
%\extrainfo{test} % Other text you want to include on this line

\position{Mechanical Engineer{\enskip\cdotp\enskip}UCSB} % Your expertise/fields
\quote{``Creativity is all about connecting dots."} % A quote or statement

%----------------------------------------------------------------------------------------
%	RECIPIENT/POSITION/LETTER INFORMATION
%	All of the below lines must be filled out
%----------------------------------------------------------------------------------------

\recipient{Internship Recruitment Team}{Lucid Motors\\Newark, CA} % The company being applied to

\letterdate{\today} % The date on the letter, default is the date of compilation

\lettertitle{Thermal Systems Intern} % The title of the letter

\letteropening{Dear Lucid Motors,} % How the letter is opened

\letterclosing{Kind regards,} % How the letter is closed

% \letterenclosure[Attached]{Resume} % Any enclosures with the letter

\makecvfooter{\today}{Jason C. Wei~~~·~~~Cover Letter}{} % Specify the letter footer with 3 arguments: (<left>, <center>, <right>), leave any of these blank if they are not needed
  
%----------------------------------------------------------------------------------------

\begin{document}

\makecvheader % Print the header

\makelettertitle % Print the title

%----------------------------------------------------------------------------------------
%	LETTER CONTENT
%----------------------------------------------------------------------------------------

\begin{cvletter}

%------------------------------------------------
As the technical director and powertrain lead for Gaucho Racing, the FSAE Electric team at UC Santa Barbara, I am delighted to submit this thermal systems intern application. I have hands-on, school sponsored research experience with thermofluidics and lab experience with corrosion resistant routing.

In the school sponsored research project, I have designed and am carrying out an experiment to characterize our FSAE vehicle's liquid cooling system. I have built a test rig with our radiator, pump, tubing, reservoir, and heating element to simulate the thermal output of the motor and motor controller. The key parameters we are trying to derive from this project is the pressure drop of our components, the heat dissipation of the radiator (radiators are very poorly characterized and generally have no datasheet), and minor losses caused by tubing geometry. Pressure, temperature, and flow rate is monitored at multiple points in the loop. This experiment will allow us to isolate component parameters and design more efficient liquid cooling systems in the future.

Lucid Air battery pack's water glycol cooling system is in the top two of my favorite Lucid innovations, the other being the mini-differential that was built within the motor. Extending the negative terminal of the battery cell to the positive side to allow for the negative side to be completely for cooling was incredibly inspiring to learn about.

At Lucid, I will bring not only my experience with thermofluidics, but also microfluidics. As a Workshop Wizard at the California NanoSystems Institute, my responsibilities include managing the microfluidics lab. I have done extensive research into using phase-change cooling (boiling of liquid) inside a heat sink with microstructures to more effectively remove heat from electronics. This is a primary research topic of our Thermofluid Energy Reserach Laboratory, which I have gotten the chance to visit and collaborate with on constructing test rigs.

I look forward to discussing this internship opportunity and my qualifications with you in further detail. Interning at Lucid would be an exciting opportunity for me to experience the industry leading technology, collaboration, and inter-team integration that happens all under the same Newark roof.

Thank you so much for your time and consideration.

%------------------------------------------------


\end{cvletter}

%----------------------------------------------------------------------------------------

\makeletterclosing % Print the signature and enclosures

\end{document}